% !PRINTED 283
\weberpage{283}

%<<CONT
formiren. Es hat sich nämlich früher schon gezeigt, dass in der Gruppe $L_7$ Theiler vom Grade 21 enthalten sind, deren Index $= 8$ ist, und die gerade für unseren Zweck von besonderer Wichtigkeit sind. Unter den conjugirten Theilern des Index 8 ist aber der einfachste und für die Rechnung bequemste der aus allen Substitutionen

\weq{5}{\begin{pmatrix} a, & b \\ 0, & a^{-1} \end{pmatrix} \quad (\mathrm{mod}\ 7)}

bestehende, und wir wollen die Transformation also so einrichten, dass $\chi$ in dieser Form auftritt. Dies erreichen wir durch Transformation der Substitutionen (4) mittelst $\begin{pmatrix} 1, & 1 \\ -1, & 0 \end{pmatrix}$, und dadurch ergiebt sich aus (4)

\weq{6}{\chi = \begin{pmatrix} 2, & 3 \\ 0, & -3 \end{pmatrix}, \ \omega = \begin{pmatrix} 1, & -3 \\ 3, & -1 \end{pmatrix}, \ \Theta = \begin{pmatrix} 1, & 3 \\ -2, & 2 \end{pmatrix}.}

Wir stellen hiernach noch zur besseren Uebersicht die ganze Octa\"edergruppe in eine Tafel zusammen, deren sechs Zeilen die Elemente $\Theta^\lambda, \omega\Theta^\lambda, \chi\Theta^\lambda, \chi\omega\Theta^\lambda, \chi^2\Theta^\lambda, \chi^2\omega\Theta^\lambda$ für $\lambda = 0, 1, 2, 3$ enthalten:

\weq{7}{\begin{gathered}
\begin{pmatrix} 1, & 0 \\ 0, & 1 \end{pmatrix}, \ \begin{pmatrix} 1, & 3 \\ -2, & 2 \end{pmatrix}, \ \begin{pmatrix} 2, & 2 \\ 1, & -2 \end{pmatrix}, \ \begin{pmatrix} 2, & -3 \\ 2, & 1 \end{pmatrix};\\
\begin{pmatrix} 1, & -3 \\ 3, & -1 \end{pmatrix}, \ \begin{pmatrix} 0, & 3 \\ 2, & 0 \end{pmatrix}, \ \begin{pmatrix} -1, & 1 \\ -2, & 1 \end{pmatrix}, \ \begin{pmatrix} 3, & 1 \\ -3, & -3 \end{pmatrix};\\
\begin{pmatrix} 2, & 3 \\ 0, & -3 \end{pmatrix}, \ \begin{pmatrix} 3, & -2 \\ -1, & 1 \end{pmatrix}, \ \begin{pmatrix} 0, & 2 \\ 3, & 1 \end{pmatrix}, \ \begin{pmatrix} 3, & -3 \\ 1, & -3 \end{pmatrix};\\
\begin{pmatrix} 3, & 2 \\ 2, & -3 \end{pmatrix}, \ \begin{pmatrix} 1, & 1 \\ -1, & 0 \end{pmatrix}, \ \begin{pmatrix} 1, & 2 \\ 1, & 3 \end{pmatrix}, \ \begin{pmatrix} -3, & 0 \\ 2, & 2 \end{pmatrix};\\
\begin{pmatrix} 3, & 3 \\ 0, & -2 \end{pmatrix}, \ \begin{pmatrix} -3, & 1 \\ -3, & 3 \end{pmatrix}, \ \begin{pmatrix} 2, & 0 \\ -2, & -3 \end{pmatrix}, \ \begin{pmatrix} -2, & 1 \\ 3, & -2 \end{pmatrix};\\
\begin{pmatrix} -2, & 2 \\ 1, & 2 \end{pmatrix}, \ \begin{pmatrix} -1, & 2 \\ 3, & 0 \end{pmatrix}, \ \begin{pmatrix} 2, & 1 \\ 3, & 2 \end{pmatrix}, \ \begin{pmatrix} 0, & 1 \\ -1, & -1 \end{pmatrix}.
\end{gathered}}
Um die ganze Gruppe $L_7$ zu erhalten, müssen wir noch ein Element $7^{\text{ten}}$ Grades hinzufügen, und dafür wählen wir

\weq{8}{\tau = \begin{pmatrix} 1, & 1 \\ 0, & 1 \end{pmatrix}.}
